\section{Introduction}

Every commutative domain embeds in a field. For noncommutative domains,
division-ring embeddability is an additional condition. We distinguish
these properties for group algebras of word-hyperbolic groups.

Recall that a group \(H\) has \emph{unique product} if, for all nonempty
finite subsets \(P,R\subset H\), some element of \(PR\) has exactly one
expression \(xy\), with \(x\in P\) and \(y\in R\). Its coefficient in a
product of two nonzero elements of \(K[H]\) cannot cancel. Thus unique
product makes \(K[H]\) a domain over every field \(K\). Our examples
satisfy the stronger condition of diffuseness, recalled in
Section~\ref{sec:diffuse}.

\begin{maintheorem}[Hyperbolic examples]\label{thm:main}
For every prime \(p\), there is a finitely presented, torsion-free,
word-hyperbolic diffuse group \(\Gamma_p\) such that
\begin{enumerate}
\item \(K[\Gamma_p]\) is a domain for every field \(K\);
\item if \(\operatorname{char}K=p\), there is no injective ring
homomorphism from \(K[\Gamma_p]\) to any division ring.
\end{enumerate}
\end{maintheorem}

Kourovka Problem~1.5, recorded by L.~A.~Bokut', asks whether there is a
group whose group ring has no zero divisors and is not embeddable into
a skew field~\cite{Kourovka}. The problem leaves the coefficient ring
unspecified; Bokut' poses the embedding question for group algebras over
fields~\cite[p.~122]{Bokut87}. Theorem~\ref{thm:main} answers it for
group algebras over fields of every positive characteristic, now with
word-hyperbolic groups. The integral version, for \(\Z[H]\), and the
version over fields of characteristic zero remain open. Mal'cev's
problem, Kourovka Problem~1.6, asks about division embeddings for
left-orderable groups and remains open as well.

\subsection*{Two constructions from one rank ceiling}

Fix a prime \(p\) and a power \(Q=p^a>320000\cdot2060^{309}\).
We construct a diffuse, torsion-free word-hyperbolic group \(G_p\)
carrying matrices \(Y_1,\ldots,Y_N\) of order \(r\), with group elements
as entries. Interpolation over \(\F_Q\) and the triangle relations give
\begin{equation}\label{eq:intro-budget}
 \sum_{i=1}^N\rk_D\phi(Y_i)<\frac{Nr}{4}
\end{equation}
for every unital homomorphism \(\phi:\F_p[G_p]\to D\) to a division
ring. The division ring need not contain \(\F_Q\): its regular
representation over \(\F_p\) supplies the required scalar matrices.
After diagonal normalization, \(Y_i\) is equivalent to
\(\diag(1,M_i)\), where
\[
 M_i=(e_{i,uv}-1)_{1\le u,v<r}
\]
and the \(e_{i,uv}\) freely generate a line subgroup \(F_i\le G_p\).

For Theorem~\ref{thm:main}, put \(n=r-1\) and take a free group
\(A=*_{i=1}^N A_i\), where \(A_i\) has basis
\(a_{i,uv}\), \(1\le u,v\le n\). For every permutation \(\pi\) of the
\(n^2\) entry positions, use a copy \(G_p^\pi\) and glue each of its
\(F_i\)'s to \(A_i\) by \(e_{i,uv}\mapsto a_{i,\pi(u,v)}\).
All vertex groups embed. In a faithful division-ring representation,
the \(a_{i,uv}-1\) are distinct and nonzero. Averaging the ranks of
all their placements gives some copy with total \(Y_i\)-rank at
least \(Nr/2\), contradicting~\eqref{eq:intro-budget}.

Convex lifts of the line graphs prove quasiconvexity and joint
malnormality. Every three-edge segment in the Bass--Serre tree has
trivial stabilizer, even though different edges at the central vertex
share the image \(A_i\). Dahmani's acylindrical combination theorem
then gives hyperbolicity. To obtain diffuseness, we prove that every
\(F_i\) is concentrated in \(G_p\). This requires a separate cone
attachment, a systole bound of \(2\), a sharper metric diffuseness
lemma, and a puncture lemma identifying the kernel as a free product.
Bowditch's graph-of-groups theorem then applies.

The shear construction remains useful for its different
forcing argument and its retraction.
\begin{maintheorem}[Retractive shear examples]\label{thm:shear}
For every prime \(p\), the base \(G_p\) is a retract of a finitely
presented diffuse group \(\Lambda_p\) such that \(K[\Lambda_p]\) is a
domain for every field \(K\) and embeds in no division ring whenever
\(\operatorname{char}K=p\). The kernel of the retraction is
\[
 *_{i,\,gF_i\in G_p/F_i}g\Z^2g^{-1}.
\]
It is bi-orderable, and \(\Lambda_p\) is left-orderable if and only if
\(G_p\) is left-orderable.
\end{maintheorem}

Here one attaches \(\Z^2\rtimes F_i\), with commuting generators
\(s_i,\gamma_i\), and relations
\[
 [s_i,e_{i,uv}]=1,\qquad
 \gamma_i e_{i,uv}\gamma_i^{-1}=s_i^{\delta_{uv}}e_{i,uv}.
\]
In the division closure of a faithful image of an attachment,
\(s_i\) is central with distinct powers. Writing
\(\Delta=\diag(e_{i,11},\ldots,e_{i,nn})\), the relations give
\begin{equation}\label{eq:intro-conjugacy}
 (\gamma_i^kI)M_i(\gamma_i^{-k}I)=\Delta(s_i^kI+C)
\end{equation}
for a fixed matrix \(C\). Singularity would give arbitrarily many
distinct central eigenvalues of \(C\), impossible in finite dimension.
Thus every \(Y_i\) has full rank, again contradicting the ceiling.
This algebraic argument has the following general form.

\begin{maintheorem}[General attachment principle]\label{thm:principle}
Let \(K\) be a field and let \(G\) be a unique-product group. Suppose
there are finitely many matrices \(Y_i\) over \(K[G]\), of orders
\(r_i\ge2\), with the following properties:
\begin{enumerate}
\item each \(Y_i\) becomes
\(\left(\begin{smallmatrix}1&\one^{\mathsf T}\\
\one&(e_{i,uv})\end{smallmatrix}\right)\)
under invertible row and column operations over \(K[G]\), where its
interior entries freely generate a subgroup \(F_i\le G\);
\item for every division ring \(D\) and injective homomorphism
\(\phi:K[G]\to D\), we have \(\sum_i\rk_D\phi(Y_i)<\sum_i r_i\).
\end{enumerate}
Then \(G\) is a retract of a unique-product group \(\Gamma\) for which
\(K[\Gamma]\) embeds in no division ring. If \(G\) is diffuse, so is
\(\Gamma\). If \(G\) is finitely presented, so is \(\Gamma\).
\end{maintheorem}

No independence condition on a division-ring embedding is assumed in
either construction. Compared with Theorem~\ref{thm:shear},
Theorem~\ref{thm:main} replaces the shear attachments, which contain
\(\Z^2\), by permutation gluing, and so yields a word-hyperbolic
example. The concentration proof for the line subgroups is what
permits that gluing while retaining diffuseness. Theorem~\ref{thm:shear}
keeps a retraction onto the base, and with it the relation to the
orderability of the base.

\subsection*{Relation to earlier work}

The marked affine lines~\cite[Lemma~2.2]{OAI252}, the triangular
relations and their link analysis~\cite[Lemma~2.3]{OAI252}, and the
interpolation identity~\cite[(2.13)--(2.14)]{OAI252} come from OpenAI's
construction of a torsion-free hyperbolic group that is not residually
finite. Douglas--Mghirbi~\cite[Theorems~1.2 and~1.4, Proposition~6.4
and Corollary~6.5]{DM} used this identity to obstruct approximate
permutation models and Linnell division rings, and separated the
vertices of the triangle complex to expose free subgroups generated by
the normalized matrix entries. We give the incidence construction,
the split complex, and the rank identity in full.

For both constructions we increase the incidence girth and change the
triangle metric to make the split group diffuse. The regular
representation of \(\F_Q\) over \(\F_p\) makes the rank argument work
for every prime \(p\). The shear argument forces full rank by central
eigenvalues; the hyperbolic argument instead forces a sufficiently
large average rank over all basis permutations.

Bowditch introduced diffuse groups and proved their unique-product,
extension and graph-of-groups properties~\cite{Bowditch}.
Our metric arguments use the strong four-point inequality of
Nica--\v Spakula~\cite[Definition~4.1 and Theorem~5.1]{NicaSpakula}.
For hyperbolic combination we use
Dahmani~\cite[Theorem~0.1(1) and Proposition~1.11]{Dahmani}; see also
Bestvina--Feighn~\cite{BF,BFAdd} and Kapovich~\cite{Kapovich}.
The broader division-embedding problem and its relation to Linnell
embeddings are discussed by Fisher--S\'anchez-Peralta~\cite{FSP}.
In 2024 Fisher and Klinge noted that no torsion-free group was known
whose group algebra fails to embed in a division ring~\cite[Section~3.2]{FK}.
The torsion-free examples found since then have group algebras with
zero divisors~\cite{DMzd,Ivanov,OAI196} or group algebras that are not
directly finite~\cite{OAI197}. Those results do not address the domain
case treated here.

\subsection*{Conventions and proof order}

Rings and homomorphisms are unital. A domain is a nonzero ring without
zero divisors. Matrix rank over a division ring is the dimension of the
right column space; row and column elimination have their usual
noncommutative meaning. A faithful representation of a group algebra
means an injective ring homomorphism.

Section~\ref{sec:rank} proves Theorem~\ref{thm:principle}.
Section~\ref{sec:lines} constructs the incidence data and proves the
interpolation identity. Section~\ref{sec:geometry} constructs \(G_p\),
its line subgroups and the rank ceiling, and proves diffuseness.
Section~\ref{sec:hyperbolic} proves Theorem~\ref{thm:main}, including
the concentration argument. Section~\ref{sec:finish} proves
Theorem~\ref{thm:shear}, specifies its presentation, and discusses the
remaining embedding and orderability questions.
